\section{Experimental Settings}
\label{sec:experiment_setting_details}

\subsection{Implementation of Stability Components}
\label{sec:implement_architecture}

To maintain numerical values within the FP16 safety margin without sacrificing model performance, we implement Logits Soft-Capping and Sandwich Normalization.
These mechanisms cap extreme values and normalize residual branches, respectively.

\paragraph{Logits Soft-Capping.}
Standard linear layers in Large Language Models often produce logits that grow unbounded during training, causing the Softmax function to saturate and gradients to vanish or explode.
Soft-capping addresses this by squashing the logits into a fixed range using the hyperbolic tangent ($\tanh$) function before scaling them back.
Formally, given the raw logits $x$ and a capping threshold $\sigma$ (e.g., $30.0$ or $50.0$), the capped logits $x'$ are computed as:
\begin{equation}
    x' = \sigma \cdot \tanh\left(\frac{x}{\sigma}\right)
\end{equation}
In our implementation, we apply this transformation to the output logits of the language model head.
This ensures that the input to the cross-entropy loss remains within the range $(-\sigma, \sigma)$, preventing logits from exceeding the FP16 maximum value while preserving the relative order of probabilities.

\paragraph{Sandwich Normalization.}
In the standard Pre-Norm Transformer architecture, the input $x$ is normalized before the sub-layer (Attention or Feed-Forward Network), and the output is added directly to the residual stream: $x_{l+1} = x_l + F(\text{Norm}(x_l))$.
While effective, this allows the magnitude of the residual stream $x$ to grow monotonically with depth, potentially destabilizing deep networks.
Sandwich Normalization introduces an additional normalization layer explicitly on the output of the sub-layer branch before the residual addition.
The modified update rule for a block containing a sub-layer $F$ (e.g., Self-Attention or MLP) is defined as:
\begin{equation}
    x_{l+1} = x_l + \text{Norm}_{\text{post}}\left( F(\text{Norm}_{\text{pre}}(x_l)) \right)
\end{equation}
In our implementation, we apply this strictly to the residual branches.
This ensures that the contribution of each layer has unit variance, preventing the accumulation of extreme activation values as the network depth increases.

\subsection{Training Configuration}

In \Cref{tab:training_config}, we present the details of our training hyperparameter configuration in three parts:
\begin{itemize}
    \item For the model architecture, we primarily follow Qwen3-1.7B~\citep{qwen3} and adopt the vocabulary from the Qwen series~\citep{qwen2,qwen3}. We use $\theta = 10000$ for RoPE~\citep{su2024rope} to support a context length of 4K. The Soft-Capping threshold is set to 30.0, as discussed in \Cref{sec:architecture_stability}.
    \item We use AdamW as the optimizer with $\beta_1 = 0.9$ and $\beta_2 = 0.95$. We adopt a $\mu$P with base dimension of 896 and set the learning rate to $5\times10^{-3}$ for Phase~1 and $3\times10^{-3}$ thereafter before decay.
    \item To support FP16 training, we use dynamic loss scaling with a factor of 2 and a window of 20 to handle widely varying gradient scales.
\end{itemize}
This detailed configuration facilitates the reproduction of our training run.

\begin{table}[htbp]
\centering
\caption{Training Hyperparameter Configuration.}
\label{tab:training_config}
\begin{tabular}{@{}llr@{}}
\toprule
\textbf{Category} & \textbf{Parameter} & \textbf{Value} \\
\midrule
\multicolumn{3}{l}{\textbf{Model Architecture}} \\
& Sequence Length & 4096 \\
& Hidden Size & 2048 \\
& FFN Dimension & 6144 \\
& Number of Layers & 28 \\
& Number of Attention Heads & 16 \\
& Number of KV Heads (GQA) & 8 \\
& Vocabulary Size & 151936 \\
& Rotary $\theta$ & 10000.0 \\
& Logit Soft-capping threshold & 30.0 \\
& Initialization Std & 0.018 \\
\midrule
\multicolumn{3}{l}{\textbf{Optimizer Configuration}} \\
& Optimizer Type & AdamW \\
& Learning Rate (Phase 1) & $5\times 10^{-3}$ \\
& Learning Rate (Phase 2+) & $3\times 10^{-3}$ \\
& Batch Size & 2048 \\
& $\beta_1$ & 0.9 \\
& $\beta_2$ & 0.95 \\
& $\epsilon$ & 1e-8 \\
& Weight Decay & 0.1 \\
& Warmup Steps & 5000 \\
& $\mu$P Width Base & 896 \\
\midrule
\multicolumn{3}{l}{\textbf{Loss Scaling (Dynamic)}} \\
& Scale Factor & 2 \\
& Scale Window & 20 \\
& Minimum Loss Scale & 524288 \\
\bottomrule
\end{tabular}
\end{table}

\subsection{Model Average}\label{sec:model_average}

\begin{table}[htbp]
\centering
\caption{Model Performance Across Checkpoints.}
\label{tab:checkpoint_performance}
\begin{narrowtblr}{c c c c c c}
Ckpt Step & ARC-Challenge & ARC-Easy & CSQA & PIQA & Average\\
260632 & 64.41 & 82.72 & 65.93 & 73.39 & 71.61\\
261032 & 65.42 & 82.19 & 65.36 & 73.72 & 71.67\\
261432 & 63.05 & 81.48 & 65.68 & 74.59 & 71.2\\
261832 & 65.42 & 81.83 & 64.78 & 73.56 & 71.40\\
262232 & 61.36 & 83.25 & 65.77 & 73.78 & 71.04\\
262632 & 65.76 & 82.01 & 66.34 & 73.5 & 71.90\\
263032 & 63.73 & 80.78 & 66.42 & 73.78 & 71.17\\
263132 & 62.71 & 80.6 & 66.09 & 73.88 & 70.82\\
\end{narrowtblr}
\end{table}

Following recent work~\citep{luo2025learningratedecaywastes}, we average the near-end checkpoints to reduce variance and consolidate learned knowledge and capabilities.
We first evaluate the last eight checkpoints on a subset of lightweight benchmarks, as shown in \Cref{tab:checkpoint_performance}. Consecutive checkpoints are spaced 400 steps apart, corresponding to 3.36B tokens.
These checkpoints fluctuate during training and do not exhibit a clear upward or downward trend. Therefore, we apply simple model averaging~\citep{li2025modelmergingpretraininglarge}, directly averaging the last eight checkpoints to obtain the final model.

\subsection{Reference Experiments for Quantile Benchmarking}
\label{sec:referece_quantile_benchmarking}

We conduct quantile benchmarking experiments across two primary scenarios: training from scratch and continual training from checkpoints. For each experiment, given a target quantile $p\%$, we select the data partition above the $p\%$ threshold, comprising roughly 10B tokens, which are then used for the respective training scenarios.

\paragraph{Training from Scratch.}
In the training-from-scratch scenario, we train a model with the Qwen3-0.6B architecture~\citep{qwen3}. Following the default configuration in \Cref{tab:training_config}, we conduct a small-scale experiment using the settings detailed in \Cref{tab:train_from_scratch}, training over approximately 8.4B tokens from the quantile data chunks. We employ a constant learning rate schedule with a sufficiently long warmup phase to ensure stable training dynamics.

\paragraph{Continual Training.}
In the continual training scenario, we resume from a checkpoint previously trained on approximately 367B tokens of the deduplicated DCLM-Baseline dataset. The model adopts the Qwen2.5-0.5B architecture~\citep{qwen2025qwen25technicalreport}. We then train over approximately 8.4B tokens from the quantile data chunks using the configuration specified in \Cref{tab:continual_training}. For these experiments, we linearly decay the learning rate from a peak value of $1\times10^{-3}$ to a final value of $1\times10^{-5}$.

\paragraph{Consistency Across Scenarios.}
As illustrated in \Cref{fig:quantile_benchmark_understanding,fig:quantile_benchmarks_knowledge}, the benchmarking results exhibit strong alignment between the training-from-scratch and continual training experiments. This consistency persists for evaluations on both the DCLM-Baseline and Fineweb-Edu datasets, despite resuming from a checkpoint trained exclusively on the deduplicated DCLM-Baseline dataset. This observation supports the robustness of our quantile-based data selection approach across different training paradigms.

\paragraph{Compute Cost Comparison.} We use the DCLM-Baseline as a reference for comparison, as discussed in \Cref{sec:data_benchmark}. Following previous work~\cite{kaplan2025scaling}, we estimate the compute budget using the formula $C=6ND$, where $N$ represents the number of model parameters and $D$ represents the data size. Consequently, the compute cost for the 0.6B model trained on 42B tokens is approximately $(0.6\times 42) / (2\times 609)\approx 2.07\%$ relative to the baseline. Similarly, we conclude that this accounts for less than 0.6\% of the total pretraining budget.


\begin{table}[htbp]
\centering
\caption{Training Hyperparameter Configuration for Quantile Benchmarking: Training from Scratch.}
\label{tab:train_from_scratch}
\begin{tabular}{@{}lr@{}}
\toprule 
\textbf{Parameter} & \textbf{Value} \\
\midrule
Learning Rate & $1\times 10^{-3}$ \\
Batch Size & 512 \\
Warmup Steps & 400 \\
Total Steps & 4000 \\
\bottomrule
\end{tabular}
\end{table}

\begin{table}[htbp]
\centering
\caption{Training Hyperparameter Configuration for Quantile Benchmarking: Continual Training.}
\label{tab:continual_training}
\begin{tabular}{@{}lr@{}}
\toprule 
\textbf{Parameter} & \textbf{Value} \\
\midrule
Peak Learning Rate & $1\times 10^{-3}$ \\
Final Learning Rate & $1\times 10^{-5}$ \\
Batch Size & 2048 \\
Total Steps & 1000 \\
\bottomrule
\end{tabular}
\end{table}

\subsection{Reference Experiments for Repetition and Curriculum Model Averaging}\label{sec:reference_repetition_curriculum}

These experiments primarily follow the experimental framework established in CMA~\citep{luo2025learningratedecaywastes}. We use a model with the Qwen2.5-1.5B architecture without tied embeddings and train on a subset of the first shard of the DCLM-Baseline dataset.

\paragraph{Baseline Configuration.}
The baseline experiment adopts uniform data ordering and employs a Warmup-Stable-Decay (WSD) learning rate schedule with a $1$-sqrt decay function~\citep{hagele,tian2025wsmdecayfreelearningrate}, decaying to a near-zero final learning rate. The detailed experimental configuration is provided in \Cref{tab:1b5_setting}.

\paragraph{High-Quality Data Utilization Strategies.}
To investigate effective high-quality data utilization, we explore two complementary approaches:
\begin{itemize}
    \item[(1)] \textbf{Repetition Strategy:} We repeat high-quality data partitions for various top-$k$ retention ratios, matching the computational FLOPs of the single-pass baseline experiment for fair performance comparison.
    \item[(2)] \textbf{Curriculum with Model Averaging:} We adopt CMA/CDMA\footnote{We do not distinguish these variants in our context, and refer to both as \textit{CMA}. By definition, the CMA method in \Cref{tab:curriculum_comparison} corresponds to the CDMA variant, which retains LR decay.}~\citep{luo2025learningratedecaywastes}, which integrates curriculum learning with either no or moderate LR decay, accompanied by model averaging over the final checkpoints.
\end{itemize}

\begin{table}[htbp]
\centering
\caption{Training Hyperparameter Configuration for Baseline and Repetition.}
\label{tab:1b5_setting}
\begin{tabular}{@{}lr@{}}
\toprule 
\textbf{Parameter} & \textbf{Value} \\
\midrule
Peak Learning Rate & $3\times 10^{-3}$ \\
Final Learning Rate & $1\times 10^{-5}$ \\
Batch Size & 512 \\
Total Steps & 15,375 \\
Decay Steps & 2,875 \\
Warmup Steps & 768 \\
\bottomrule
\end{tabular}
\end{table}

\begin{table}[htbp]
\centering
\caption{Training Hyperparameter Configuration for Curriculum Model Average.}
\label{tab:1b5_curriculum_setting}
\begin{tabular}{@{}lr@{}}
\toprule 
\textbf{Parameter} & \textbf{Value} \\
\midrule
Peak Learning Rate & $3\times 10^{-3}$ \\
Final Learning Rate & $1\times 10^{-3}$ \\
Batch Size & 512 \\
Total Steps & 15,375 \\
Decay Steps & 2,875 \\
Warmup Steps & 768 \\
Checkpoint Number & 6\\
Decay Factor of EMA ($\alpha$) & 0.2\\
Checkpoint Interval & 0.21B\\
\bottomrule
\end{tabular}
\end{table}

\paragraph{Experimental Variants.}
The repetition experiments follow identical settings to the baseline, differing only in dataset construction. For the curriculum experiments, we use a higher final learning rate of $1\times10^{-3}$ and perform an exponential moving average (EMA) over the final six checkpoints (the last-step checkpoint is weighted by $(1-\alpha)$ relative to the current-step checkpoint, where $\alpha$ is the decay factor), replicating the methodology from CMA~\citep{luo2025learningratedecaywastes}.

\paragraph{Evaluation Settings.}
In \Cref{tab:curriculum_comparison}, we evaluate performance on a high–signal-to–noise-ratio benchmark subset (\textit{Core} in \Cref{tab:curriculum_comparison}) comprising MMLU~\citep{mmlu}, ARC~\citep{arc}, and CSQA~\citep{commonsenseqa}, following established practices in prior work~\citep{heineman2025signalnoiseframeworkreducing,luo2025learningratedecaywastes}. These benchmarks provide strong discriminative power for identifying performance differences between training approaches.
























